\documentclass[12pt]{article}
\usepackage[utf8]{inputenc}
\usepackage[T1]{fontenc}
\usepackage{hyperref}
\title{Lepontic Language Resources for K-12 Curriculum Development}
\author{Compiled by Claude Sonnet 5.5 (claude-sonnet-5-5) for Andr\'as Kornai}
\date{}
\begin{document}
\maketitle

\section{Introduction}
Lepontic is an extinct Celtic language attested in about 150 short inscriptions from the Lugano area, the north Italian lake region and the Swiss canton of Ticino, dating from roughly the 6th to the 1st century BC. They are the earliest written evidence of any Celtic language. The texts are in the Lugano (Lepontic) alphabet, a North Italic script derived from Etruscan. Almost only names are attested, and the lexicon is small. Scholars disagree whether Lepontic is a distinct Continental Celtic language or an early form of Gaulish, and the boundary with Cisalpine Gaulish is hard to draw. There are no native speakers, and no school materials were found.

\section{Primary Educational Resources}
No primers, textbooks or school programmes were found in search results.

\section{Linguistic and Reference Materials}
\begin{itemize}
\item LexLep, the Lexicon Leponticum, a digital edition of the Cisalpine Celtic inscriptions: ``The Cisalpine Celtic Languages''. \url{https://lexlep.univie.ac.at/wiki/The_Cisalpine_Celtic_Languages}
\item Wikipedia, ``Lepontic language''. \url{https://en.wikipedia.org/wiki/Lepontic_language}
\item Omniglot, Lepontic writing: a description of the script. \url{https://omniglot.com/writing/lepontic.htm}
\item Stifter, ``Lepontic'' (scribd copy; I did not verify its content or rights status). \url{https://es.scribd.com/document/65123135/Stifter-Lepontic}
\item Humboldt University repository: ``Lepontische Inschriften aus dem Brembanatal''. \url{https://edoc.hu-berlin.de/18452/34440}
\item Roots of Europe (University of Copenhagen): ``Lepontisch'' course document (2010). \url{https://rootsofeurope.ku.dk/kalender/arkiv_2012/celtic_spring/lepontisch_ws_2010.pdf}
\end{itemize}

\section{Religious and Community Resources}
None were found.

\section{Digital and Audio Resources}
The LexLep project (above) is the main digital resource. No audio resources or software were found.

\section{Curriculum Development Context}
A heritage unit could use the inscriptions, the Lugano alphabet and the Celtic connection, taking examples from LexLep. It would need to state plainly that the attested vocabulary is mostly names and that the language's classification is disputed.

\end{document}